\section{Detection and harmonic signatures}\label{supp:detection}
These calculations support the carrier-scaling and harmonic tests in
\S\ref{sec:transduction} and the dynamic detector in \S\ref{sec:dynamic}.
\subsection{Fixed-voltage mixing and carrier exponents}\label{supp:dynamic-details}
Time and frequency use the membrane-time units defined in article
Appendix~\ref{app:mpr-coefficients}; $r^*,d,\nu$ have the same definitions.
Membrane access is already included in this calculation. The linear voltage response
is $H_v(\omega)=(i\omega+d)/[(i\omega+d)^2+\nu^2]$. Dividing by its squared magnitude
gives the mixing coefficient at fixed voltage excursion,
\begin{equation}
\kappa_V(Q,\omega)=\frac{\kappa_{\rm DC}(Q,\omega)}{|H_v(\omega)|^2}
=\frac{F'(Q)}{2}\;\frac{\omega^2+3d^2-\nu^2}{\omega^2+d^2}.
\label{eq:kfixed}
\end{equation}
This coefficient tends to $F'/2$ at high frequency. The $1/\omega^2$ suppression
at fixed current comes from $|H_v|^2$, the membrane filtering already represented
by \ifsupp main-text\fi\ Eq.~\eqref{eq:doors}; it is not an additional detector filter. A static curve
placed after the same membrane has the same leading carrier roll-off, although
its mixing coefficient and sign can differ. The steady-state curvature is
$F''=4r^*(3d^2-\nu^2)/(d^2+\nu^2)^3$, so the inflection occurs at $\nu^2=3d^2$.

For nearby fast tones, the quadratic terms in \ifsupp main-text\fi\ Eq.~\eqref{eq:mpr} contribute at
different orders: $v^2$ at $\omega^{-2}$, $r v$ at $\omega^{-3}$, and $r^2$ at
$\omega^{-4}$. Table~\ref{tab:exponents} compares the carrier and beat routes to
the firing rate for a static curve and a dynamic population with the same membrane access.

\begin{table}[h]
\centering\small
\begin{tabular}{lcc}
\toprule
stage & static curve, same membrane & dynamic population \\
\midrule
point-membrane current $\to$ voltage & $\omega^{-1}$ & $\omega^{-1}$ \\
voltage $\to$ rate, carrier & $1$ & $\omega^{-1}$ \\
current $\to$ rate, carrier & $\omega^{-1}$ & $\omega^{-2}$ \\
two tones $\to$ beat in the rate & $(\omega_1\omega_2)^{-1}$ & $(\omega_1\omega_2)^{-1}\times$ const. \\
\bottomrule
\end{tabular}
\caption{Fast-carrier scaling with matched point-membrane access, at fixed beat frequency
and a nonzero leading mixing coefficient. The dynamic population integrates the
carrier once more between voltage and rate; the beat is formed before that step.
Prefactors depend on the operating point and beat-frequency susceptibility.}
\label{tab:exponents}
\end{table}

The first row gives the shared membrane factor $1/\omega$. The second converts
voltage to rate: a static curve adds its slope, whereas the dynamic rate equation
integrates the carrier once more. Combining these stages gives the third-row
carrier responses, $\omega^{-1}$ and $\omega^{-2}$, respectively.

The beat in the last row is formed before that additional rate integration:
static curvature and dynamic $v^2$ each mix the attenuated voltage, retaining
$(\omega_1\omega_2)^{-1}$. Equal exponents do not imply equal amplitudes; mixing
and the susceptibility at fixed $\Omega$ determine the prefactors. For the
uncoupled carrier-averaged response, $\kappa_V$ in Eq.~\eqref{eq:kfixed} tends to
$F'/2$, but is not the full finite-$\Omega$ network prefactor. Thus the dynamic
rate attenuates the carrier twice while forming the beat between those stages.

\ifjne
\subsection{Harmonic ladder and mixed detectors}\label{supp:harmonics}
Envelope following and square-law detection are distinguishable at a single site. At a matched target,
$\mathcal A(t)=2\varepsilon|\cos(\Omega t/2)|$, whose Fourier series
$|\cos\theta|=\tfrac2\pi+\tfrac4\pi\sum_{n\ge1}\tfrac{(-1)^{n+1}}{4n^2-1}\cos(2n\theta)$ gives a
harmonic ladder whose nonnegative Fourier amplitudes $A_{n\Omega}$ obey $A_{2\Omega}/A_\Omega=1/5$, $A_{3\Omega}/A_\Omega=3/35$ and
$A_{4\Omega}/A_\Omega=1/21$. The square-law output $\half \mathcal A^2=\varepsilon^2(1+\cos\Omega
t)$ contains only DC and $\Omega$. For the mixed family $T(\mathcal A)=c_1\mathcal A+c_2\mathcal A^2$, with
$c_1\ne0$ and a nonzero fundamental, the
pre-network ratio is
\begin{equation}
\mathcal H_{\rm det}\equiv\left|\frac{Y_{2\Omega}}{Y_\Omega}\right|
=\frac{1/5}{\left|1+\tfrac{3\pi}{4}\,c_2\varepsilon/c_1\right|},
\label{eq:rdet}
\end{equation}
so measuring $\mathcal H_{\rm det}$ constrains $c_2\varepsilon/c_1$ within this family.\footnote{The inverse has two branches unless the sign of the fundamental
relative to the linear term is known; for $c_2/c_1\ge0$ it is unique.}

\fi
\subsection{AM-surrogate and genuine two-tone comparisons}\label{supp:waveform-checks}
These checks support the waveform choices in \S\ref{sec:waveforms}. Agreement
at the leading beat does not imply equal steady sources or driven operating states.

\begin{table}[htbp]\centering\small
\caption{Waveform comparisons at the stated operating points. JR response ratios
use the same voltage readout; timing ratios use the same pyramidal firing rate.}
\label{tab:waveform-checks}
\begin{tabular}{@{}p{4.3cm}p{4.5cm}p{6cm}@{}}
\toprule
Comparison & Conditions & Result\\
\midrule
JR open-loop beat & $v^*=8$ mV & Two-tone/AM: $1.0000$ at $0.02$ mV; $0.988$ at $1$ mV.\\
JR modulation/mean shift & Pinned sweep: far from Hopf, $\gamma=1$ s$^{-1}$, nearest point & Two tones: $14.5$, $58$, $44$; AM: $19$, $66$, $45$. AM has $3/4$ of the two-tone direct steady source.\\
JR TI/tACS output & Far from Hopf, $\gamma=1.07$ and $0.12$ s$^{-1}$ & AM: $0.14$, $0.24$, $0.47$; two tones: $0.14$, $0.24$, $0.46$. Difference at most $2.2\%$ across the grid.\\
JR matched-source output & Three couplings, settings in \S\ref{app:jr-sweep-settings} & TI/tACS: $0.99$, $0.94$, $0.92$ for AM; $0.99$, $0.93$, $0.90$ for two tones.\\
PING beat & Well below Hopf; $\bar\eta=0.9$ & Matched AM and two-tone responses agree within $0.07\%$ and $4\%$, respectively.\\
\bottomrule
\end{tabular}
\end{table}

\subsection{The cubic Hopf normal form lacks a quadratic mixer}\label{app:sl}
The Stuart--Landau equation \cite{pikovsky2001}
\begin{equation}
\dot z=(\alpha+i\omega_0)z-|z|^2z+\mathcal F(t),\qquad
\gamma=-\alpha>0,
\label{eq:sl}
\end{equation}
uses seconds for time and a normalized complex mode $z$, with the cubic
coefficient set to $1$~s$^{-1}$. Its linear growth rate is $\alpha=-\gamma$;
$\gamma$ has units s$^{-1}$. The mode has damping $\gamma$ and natural angular frequency
$\omega_0$, driven by the applied input $\mathcal F$ (normalized mode units per second). For a weak input
$\mathcal F=\varepsilon_{\rm SL} \psi_{\rm SL}(t)$ with dimensionless waveform $\psi_{\rm SL}$ on the stable side, the perturbative response has
the form $z=\varepsilon_{\rm SL} z_1+\varepsilon_{\rm SL}^3z_3+\cdots$. There is no second-order
source: the first-order response contains only the applied frequencies, and the
cubic first contributes at order $\varepsilon_{\rm SL}^3$. Thus the normal form has a
resonant linear susceptibility but lacks the quadratic difference-frequency
generation of \ifsupp main-text\fi\ Eq.~\eqref{eq:demod}. This statement concerns the weak-input
expansion; it is not an exclusion of every higher-order intermodulation or
strong-drive response.

In the numerical control ($\alpha=-0.5$~s$^{-1}$, $f_0=10$~Hz, $f_c=100$~Hz,
$\varepsilon_{\rm SL}=0.5$~s$^{-1}$), a square-law detector gives a beat-frequency lock-in amplitude
of $2.5\times10^{-1}$, while $\mathrm{Re}\,z$ gives approximately
$2\times10^{-8}$, at the numerical floor. Driving the same mode directly in band
recovers the forced $1/\gamma$ growth. The control separates resonant amplification
from the nonlinear source required to drive it at the beat.
